\documentclass[aspectratio=169,12pt]{beamer}
\usepackage{amsmath}
\usepackage{amssymb}
\usepackage{array}
\usepackage[most]{tcolorbox}
\usepackage{graphicx}
\usepackage{tikz}
\usepackage[sfdefault,lf]{FiraSans}
\renewcommand{\familydefault}{\sfdefault}
\setlength{\parindent}{0pt}
\everymath{\displaystyle}
\setbeamertemplate{navigation symbols}{}
\setbeamertemplate{footline}{}
\setbeamertemplate{headline}{}
\definecolor{mmpageWhite}{HTML}{FCFBF7}
\definecolor{mmtanSoft}{HTML}{F7F5EF}
\definecolor{mmgreenDeep}{HTML}{244B2C}
\definecolor{mmgreenLeaf}{HTML}{5D8A56}
\definecolor{mmgreenSoft}{HTML}{E4EFD9}
\definecolor{mmtext}{HTML}{163221}
\definecolor{mmwarn}{HTML}{8F2F36}
\definecolor{mmwarnSoft}{HTML}{F8E8EA}
\setbeamercolor{background canvas}{bg=mmpageWhite}
\setbeamercolor{normal text}{fg=mmtext,bg=mmpageWhite}
\setbeamercolor{itemize item}{fg=mmgreenDeep}
\newtcolorbox{MMCard}[2][]{enhanced,breakable=false,colback=#2,colframe=black,boxrule=1.5pt,arc=8pt,left=5pt,right=5pt,top=4pt,bottom=4pt,before skip=0pt,after skip=0pt,#1}
\newcommand{\KoruIcon}{\raisebox{-0.2em}{\includegraphics[height=1.05em]{../../../manamaths/SITE/header-logo.png}}}
\newcommand{\NotesTitle}[1]{{\colorbox{mmtanSoft}{\parbox{0.975\linewidth}{\vspace{0.14em}\hspace{0.25em}{\LARGE\bfseries\textcolor{mmgreenDeep}{#1}\hfill\KoruIcon}\vspace{0.14em}}}\par\vspace{0.18em}{\color{mmgreenLeaf}\rule{\linewidth}{2.0pt}}\par\vspace{0.12em}}}
\newcommand{\SectionTitle}[1]{{\large\bfseries\textcolor{mmgreenDeep}{#1}}\par\vspace{0.04em}}
\newcommand{\GreenDivider}{\par\vspace{0.01em}{\color{mmgreenLeaf}\rule{\linewidth}{2.4pt}}\par\vspace{0.03em}}
\newcommand{\KeyIdea}[1]{\begin{MMCard}[title={\textbf{Key idea}},colbacktitle=mmgreenSoft,coltitle=mmgreenDeep]{mmtanSoft}\small #1\end{MMCard}}
\newcommand{\WorkedExample}[2]{\begin{MMCard}[title={\textbf{#1}},colbacktitle=mmgreenSoft,coltitle=mmgreenDeep]{white}\small #2\end{MMCard}}
\newcommand{\CommonMistake}[1]{\begin{MMCard}[title={\textbf{Common mistake}},colbacktitle=mmwarnSoft,coltitle=mmwarn]{white}\small #1\end{MMCard}}
\newcommand{\TryThese}[1]{\begin{MMCard}[title={\textbf{Try these}},colbacktitle=mmgreenSoft,coltitle=mmgreenDeep]{mmtanSoft}\small #1\end{MMCard}}
\newcommand{\StepLine}[2]{\textbf{#1.} #2\par\vspace{0.03em}}
\setbeamertemplate{background}{\begin{tikzpicture}[remember picture,overlay]\draw[black,line width=1.5pt,rounded corners=10pt]([xshift=0.45cm,yshift=-0.45cm]current page.north west) rectangle ([xshift=-0.45cm,yshift=0.45cm]current page.south east);\end{tikzpicture}}
\begin{document}
\begin{frame}[t,shrink=10]
\NotesTitle{Notes \& Steps}
\footnotesize
\begin{columns}[T,onlytextwidth]
\begin{column}{0.48\textwidth}
\KeyIdea{An angle is the amount of turn between two rays from a common vertex. Angles are measured in \textbf{degrees (°)} using a \textbf{protractor}.}
\SectionTitle{How to use a protractor}
\StepLine{1}{Place the centre hole or crosshair exactly on the \textbf{vertex} (the corner).}
\StepLine{2}{Line up the 0° mark with one arm of the angle.}
\StepLine{3}{Read the number where the other arm crosses the scale.}
\StepLine{4}{Check which scale (inner or outer) starts at 0° and use that one.}
\vspace{0.15em}
\KeyIdea{Angle types: \\
\textbf{Acute} = less than 90° \\\
\textbf{Right} = exactly 90° \\\
\textbf{Obtuse} = between 90° and 180° \\\
\textbf{Straight} = exactly 180° \\\
\textbf{Reflex} = more than 180°}
\end{column}
\begin{column}{0.48\textwidth}
\SectionTitle{Protractor tips}
\begin{itemize}
\item Always start from 0°, not from the edge of the protractor
\item If the angle opens to the left, use the outer scale; if to the right, use the inner scale (or vice versa — check which one starts at 0)
\item For accuracy, estimate the angle first — if it looks like 30°, don't read 150°
\item The two scales add up to 180° (inner + outer = 180°)
\end{itemize}
\vspace{0.2em}
\CommonMistake{Reading the wrong scale. Estimate the angle type first. If the angle is acute (looks small), the answer should be less than 90°. If you read 120° on an acute-looking angle, you probably used the wrong scale.}
\end{column}
\end{columns}
\end{frame}
\begin{frame}[t,shrink=12]
\NotesTitle{Notes \& Steps}
\begin{columns}[T,onlytextwidth]
\begin{column}{0.48\textwidth}
\SectionTitle{Estimating angles}
Before measuring, practise estimating:
\begin{itemize}
\item A full turn = 360°
\item Half a turn = 180° (straight line)
\item Quarter turn = 90° (right angle)
\item A slightly open door ≈ 30°
\item A wide open door ≈ 120°
\end{itemize}
\WorkedExample{Example 1}{Estimate then measure: the angle looks about half a right angle (≈ 45°). When measured, it reads 42° — close!}
\WorkedExample{Example 2}{Draw a 60° angle: draw a baseline ray, place protractor at endpoint, mark at 60°, draw second ray through the mark.}
\end{column}
\begin{column}{0.48\textwidth}
\TryThese{\begin{enumerate}
  \item Estimate: what type of angle is 85°?
  \item Estimate: what type of angle is 195°?
  \item If the inner scale reads 40°, what does the outer scale read?
  \item A student reads 130° for an acute-looking angle. What went wrong?
\end{enumerate}}
\CommonMistake{Not placing the protractor centre on the vertex. Even being 1 mm off can cause a 2--3° error. The measurement must start from the exact vertex.}
\CommonMistake{Using the wrong baseline. The 0° mark must align with one ray. If the protractor is tilted, all readings will be wrong.}
\end{column}
\end{columns}
\end{frame}
\end{document}
